\section{Results}
\label{sec:results}

\subsection{Main DiD Event Study}

Table~\ref{tab:main_did} presents both the continuous and binary treatment specifications with ISCO-clustered standard errors (40 clusters).

The continuous specification (left panel) yields negative interaction coefficients, indicating that more teleworkable occupations experienced smaller informality increases. A one-unit increase in the teleworkability score is associated with a 19--22 pp smaller informality increase: $\gamma_{2021} = -0.202$ ($p = 0.092$), $\gamma_{2022} = -0.220$ ($p = 0.049$), $\gamma_{2023} = -0.188$ ($p = 0.080$), $\gamma_{2024} = -0.190$ ($p = 0.046$). The year effects ($\beta_k \approx 0.17$) capture the common informality increase for workers at teleworkability = 0.

The binary specification (right panel) shows that contact-intensive workers (TW $< 0.20$, 74.6\% of the sample) experienced a differential informality increase of 9.0 pp in 2021, though this is not statistically significant at the 5\% level ($p = 0.153$). By 2024 the differential is 8.8 pp ($p = 0.062$). The weaker significance in the binary specification reflects both the coarse treatment definition (grouping 74.6\% of workers as ``treated'') and the conservative ISCO-clustered standard errors.

\subsection{Persistence Test}

The Wald test of $H_0: \gamma_{2024} = \gamma_{2021}$ yields $p = 0.937$ for the continuous specification, failing to reject the null. The interaction coefficient in 2024 ($-0.190$) is virtually identical to the 2021 value ($-0.202$), indicating that the population-level teleworkability gap in informality has not narrowed over four years. We interpret this as evidence of persistence in the differential, while noting that this does not establish within-individual permanence given the rotating panel design.

\subsection{Within-Estimator Results}

The individual fixed effects specification yields somewhat larger interaction coefficients: $-0.256$ in 2021 declining modestly to $-0.237$ in 2024, all significant at the 5\% level. The slight attenuation may reflect compositional changes in the rotating panel or genuine within-individual partial recovery for workers observed in consecutive years.

\subsection{Heterogeneity}

Figure~\ref{fig:heterogeneous} presents heterogeneous effects estimated separately by subgroup using the binary treatment within-estimator. We report the 2024 interaction coefficients:

\paragraph{Gender.} Men in contact-intensive occupations show a differential of 5.0 pp ($p = 0.224$), while women show a larger differential of 11.4 pp ($p = 0.114$). Neither is individually significant at conventional levels, but the larger female coefficient is consistent with women facing greater barriers to re-formalization in contact-intensive sectors.

\paragraph{Lima versus Rest of Peru.} The differential is significant in Lima (7.3 pp, $p = 0.001$) but not in the rest of Peru (8.1 pp, $p = 0.175$). The precision in Lima likely reflects its more homogeneous labor market and lower within-cluster variance.

\paragraph{Age Groups.} Prime-age workers (30--44) show the largest differential (8.7 pp, $p = 0.074$), followed by seniors 45--65 (8.5 pp, $p = 0.166$) and youth 15--29 (6.9 pp, $p = 0.283$). The weaker youth effect may reflect higher baseline informality among young workers, creating a ceiling effect.
